\section{Methodology}
\label{sec:method}

\subsection{Dataset Construction}
\label{sec:data}

We construct a dataset of English words with IPA transcriptions by intersecting the \wikipron US English broad transcription database \citep{lee2020massively} with the set of words that encode as a single token under the GPT-2 tokenizer.
This single-token constraint ensures a clean mapping between token representations and phonemic content: each word's embedding directly corresponds to its complete pronunciation without BPE segmentation artifacts.

The resulting dataset contains 13{,}105 words spanning 77 unique IPA phonemes.
Of these, 4{,}607 words (35\%) appear in the \oxfordfive common word list, confirming broad coverage of everyday English.
We use the first transcription listed in \wikipron for words with multiple pronunciation entries.
Each word is labeled with a multi-hot vector indicating which IPA phonemes it contains (\eg ``water'' maps to the set of its constituent IPA phonemes).
We split the data 80/20 into train (10{,}484) and test (2{,}621) sets with random seed~42.

\subsection{Experiment 1: Layer-Wise Phoneme Probing}
\label{sec:method_probing}

{\bf Where does phonetic information live across layers?}
For each of the 25 representation sites in \gptmed---the token embedding layer plus 24 transformer layer outputs---we extract the residual stream activation $\vh_\ell \in \sR^{1024}$ for each word.
We extract activations using TransformerLens \citep{nanda2022transformerlens}, processing words in batches of 256 on a single NVIDIA RTX A6000 GPU.

At each layer $\ell$, we train 77 independent Ridge classifiers (one per IPA phoneme $p$), each predicting whether phoneme $p$ is present in the word from the activation vector $\vh_\ell$.
We use Ridge regression ($\alpha = 1.0$) with StandardScaler normalization rather than logistic regression: preliminary experiments showed logistic regression required approximately 10 minutes per layer, while Ridge achieved comparable quality in approximately 11 seconds per layer, enabling a full 25-layer sweep.

We evaluate with macro F1 (averaging across all 77 phonemes equally) and micro F1 (weighting by phoneme frequency).
Macro F1 is our primary metric as it better captures performance on rare phonemes.

\subsection{Experiment 2: Phoneme Subspace Geometry}
\label{sec:method_geometry}

{\bf Does each IPA sound have a distinct direction in representation space?}
For each phoneme $p$ with at least 30 occurrences, we compute its \emph{direction vector} as the difference between the mean activation of words containing $p$ and words not containing $p$:
\begin{equation}
    \vd_p = \frac{1}{|S_p|}\sum_{\vx \in S_p} \vx - \frac{1}{|S_p^c|}\sum_{\vx \in S_p^c} \vx
    \label{eq:direction}
\end{equation}
where $S_p = \{\vx : p \in \text{phonemes}(\vx)\}$ is the set of activations for words containing phoneme $p$.
We compute these directions at every layer and analyze them via:

\begin{itemize}[leftmargin=*,itemsep=2pt,topsep=2pt]
    \item \textbf{Pairwise cosine similarity.} We compute the $49 \times 49$ matrix of absolute cosine similarities between all phoneme direction vectors and compare the mean to the expected value for random unit vectors in $\sR^{1024}$ ($\approx 0.025$).
    \item \textbf{Clustering analysis.} We apply hierarchical clustering (Ward's method) to the direction vectors and compute silhouette scores for groupings by vowel/consonant distinction and by manner of articulation.
    \item \textbf{PCA visualization.} We project all phoneme directions, vowels only, and consonants only onto their top principal components.
    \item \textbf{Cross-layer comparison.} We track how the mean pairwise similarity and clustering quality evolve from the embedding layer through the final transformer layer.
\end{itemize}

We use direction vectors rather than Ridge regression weights because the probe weights exhibited near-collinear structure---an artifact of L2 regularization---while mean-difference directions directly measure how each phoneme is encoded in the representation space.

\subsection{Experiment 3: Context-Dependent Rhyming}
\label{sec:method_context}

{\bf Can models resolve pronunciation when it depends on context?}
We test whether rhyming knowledge extends beyond static lexical representations by querying \gptfour on heteronyms: words whose pronunciation changes with context (\eg ``read'' as /ri:d/ in ``I will read'' vs.\ /r$\varepsilon$d/ in ``I have read'').

We select 12 common English heteronyms (read, lead, wind, tear, bow, bass, dove, minute, close, live, refuse, desert) and construct sentence pairs that establish each pronunciation through syntactic or semantic context.
For each context, we present a forced-choice rhyming question with two candidate words (one matching each pronunciation) and record the model's selection.
We also include 20 standard (non-heteronym) rhyming pairs as controls: 15 easy pairs and 5 orthographically tricky pairs (\eg ``cough'' and ``off'').
All queries use temperature~0 for deterministic outputs.
